\section{Definitions}

\bigskip

\begin{definition}
\label{def:Fin}
\href{https://leanprover-community.github.io/mathlib4_docs/Init/Prelude.html#Fin}{$\Fin$}$ n :=\{0,\ldots,n-1\}$, $n\in\mathbb{N}$
($\Fin 0 := \emptyset$). $\Zz := \{0,1\}$ is the \href{https://leanprover-community.github.io/mathlib4_docs/Mathlib/Algebra/Ring/Defs.html#CommRing}{commutative ring} of \href{https://leanprover-community.github.io/mathlib4_docs/Mathlib/Data/ZMod/Defs.html#ZMod}{integers modulo} $2$.
\end{definition}

\medskip

\begin{definition}
\label{def:Zn}
$\Zz^n := \Zz\times...\times\Zz$. $\Zz^n$ can be seen as the set of all functions from $\Fin n$ to $\Zz$. In this file, we consider $\Zz^n$ as a \href{https://leanprover-community.github.io/mathlib4_docs/Mathlib/Algebra/Module/Defs.html#Module}{module} over $\Zz$.
\end{definition}

\medskip

\begin{definition}
\label{def:e}
For $i\in \Fin n$, let ${\rm e}_i\in\Zz^n$ be the canonical vector with
$1$ on the $i$-th position and zeros elsewhere
\begin{equation}
\tag{e}
\label{e}
{\rm e}_i := (0,\ldots,0,1,0,\ldots,0),
\end{equation}
${\rm e}_i$ can be also seen as the \href{https://leanprover-community.github.io/mathlib4_docs/Mathlib/Algebra/Notation/Pi/Basic.html#Pi.single}{function defined on $\Fin n$,
supported at $i$, with value $1$ there, and $0$ elsewhere.}


Note that $b+{\rm e}_i$ is the code $b$ with the $i$-th bit flipped,
for every $b\in\Zz^n$ and $i\in \Fin n$.
\end{definition}

\medskip

\begin{definition}
\label{def:SF}
We say that $f:\Zz^n \to \Fin n$ is a {\bf strategy function}, if
\begin{equation}
\label{SF}
\tag{SF}
\forall\, b\in\Zz^n, k\in\Fin n, \exists\, i \in \Fin n, f(b+{\rm e}_i) = k.
\end{equation}

We say that the puzzle {\bf has a strategy} for $n$ if
$$\exists f: \Zz^n \to \Fin n, \; \eqref{SF}.$$
\end{definition}

\medskip

\begin{definition}
\label{def:d}
Let $\delta: \Fin n \times \Fin n \to \mathbb{N}$ be given by
\begin{equation}
\tag{$\delta$}
\label{d}
\delta(i,j) \; = \; \B \text{if } i=j \text{ then } 1 \text{ else } 0.
\end{equation}
See the Mathlib definition \href{https://leanprover-community.github.io/mathlib4_docs/Init/Prelude.html#if-then-else}{if-then-else}.
\end{definition}
